% Options for packages loaded elsewhere
\PassOptionsToPackage{unicode}{hyperref}
\PassOptionsToPackage{hyphens}{url}
\PassOptionsToPackage{dvipsnames,svgnames,x11names}{xcolor}
\documentclass[
  english,
  10pt,
  a4paper,
]{article}
\usepackage{xcolor}
\usepackage[top=22mm,bottom=20mm,left=18mm,right=18mm]{geometry}
\usepackage{amsmath,amssymb}
\setcounter{secnumdepth}{-\maxdimen} % remove section numbering
\usepackage{iftex}
\ifPDFTeX
  \usepackage[T1]{fontenc}
  \usepackage[utf8]{inputenc}
  \usepackage{textcomp} % provide euro and other symbols
\else % if luatex or xetex
  \usepackage{unicode-math} % this also loads fontspec
  \defaultfontfeatures{Scale=MatchLowercase}
  \defaultfontfeatures[\rmfamily]{Ligatures=TeX,Scale=1}
\fi
\usepackage{lmodern}
\ifPDFTeX\else
  % xetex/luatex font selection
\fi
% Use upquote if available, for straight quotes in verbatim environments
\IfFileExists{upquote.sty}{\usepackage{upquote}}{}
\IfFileExists{microtype.sty}{% use microtype if available
  \usepackage[]{microtype}
  \UseMicrotypeSet[protrusion]{basicmath} % disable protrusion for tt fonts
}{}
\makeatletter
\@ifundefined{KOMAClassName}{% if non-KOMA class
  \IfFileExists{parskip.sty}{%
    \usepackage{parskip}
  }{% else
    \setlength{\parindent}{0pt}
    \setlength{\parskip}{6pt plus 2pt minus 1pt}}
}{% if KOMA class
  \KOMAoptions{parskip=half}}
\makeatother
\usepackage{longtable,booktabs,array}
\usepackage{caption}
\captionsetup[table]{skip=6pt}
\newcounter{none} % for unnumbered tables
\usepackage{calc} % for calculating minipage widths
% Correct order of tables after \paragraph or \subparagraph
\usepackage{etoolbox}
\makeatletter
\patchcmd\longtable{\par}{\if@noskipsec\mbox{}\fi\par}{}{}
\makeatother
\ifLuaTeX
\usepackage[bidi=basic,shorthands=off]{babel}
\else
\usepackage[bidi=default,shorthands=off]{babel}
\fi
\ifLuaTeX
  \usepackage{selnolig} % disable illegal ligatures
\fi
\setlength{\emergencystretch}{3em} % prevent overfull lines
\providecommand{\tightlist}{%
  \setlength{\itemsep}{0pt}\setlength{\parskip}{0pt}}
% ---- Çılgın Kurye document theme: same look as the CSE 396 project proposal ----
% Compile with XeLaTeX (also works on Overleaf). Fonts: TeX Gyre Termes (Times), DejaVu Sans Mono.
\usepackage{fontspec}
\setmainfont{texgyretermes}[Extension=.otf, UprightFont=*-regular, BoldFont=*-bold, ItalicFont=*-italic, BoldItalicFont=*-bolditalic]
\setsansfont{texgyreheros}[Extension=.otf, UprightFont=*-regular, BoldFont=*-bold, ItalicFont=*-italic, BoldItalicFont=*-bolditalic]
\setmonofont{DejaVuSansMono}[Extension=.ttf, UprightFont=*, BoldFont=*-Bold, ItalicFont=*-Oblique, BoldItalicFont=*-BoldOblique, Scale=0.80]
\setmathfont{texgyretermes-math.otf}

\usepackage{xcolor}
\definecolor{ckblue}{HTML}{1F3A5F}
\definecolor{cklight}{HTML}{E8EEF5}
\definecolor{ckgrey}{HTML}{F2F2F2}
\definecolor{ckline}{HTML}{555555}
\definecolor{shadecolor}{HTML}{F4F5F7}

% Headings as in the proposal: serif, bold; subsections bold italic. Numbers are part of the heading text.
\setcounter{secnumdepth}{0}
\usepackage{titlesec}
\titleformat{\section}{\rmfamily\Large\bfseries}{}{0pt}{}
\titleformat{\subsection}{\rmfamily\large\bfseries\itshape}{}{0pt}{}
\titleformat{\subsubsection}{\rmfamily\normalsize\bfseries}{}{0pt}{}
\titlespacing*{\section}{0pt}{12pt}{4pt}
\titlespacing*{\subsection}{0pt}{8pt}{3pt}

% Header and footer (the first page has only the footer, like the proposal)
\usepackage{fancyhdr,lastpage}
\pagestyle{fancy}
\fancyhf{}
\fancyhead[L]{\footnotesize\itshape CSE 396 · Group 8 · Çılgın Kurye}
\fancyhead[R]{\footnotesize\itshape\ckmodule}
\fancyfoot[C]{\footnotesize Page \thepage\ of \pageref*{LastPage}}
\renewcommand{\headrulewidth}{0.4pt}
\setlength{\headheight}{14pt}
\fancypagestyle{ckfirst}{\fancyhf{}\renewcommand{\headrulewidth}{0pt}\fancyfoot[C]{\footnotesize Page \thepage\ of \pageref*{LastPage}}}

% Inline code: breaks only after . : ( , / _ and spaces, never inside a name
\usepackage[obeyspaces,spaces]{url}
\AtBeginDocument{%
  \renewcommand{\UrlBreaks}{\do\.\do\:\do\(\do\,\do\/\do\_\do\|}%
  \renewcommand{\UrlBigBreaks}{}%
  \renewcommand{\UrlNoBreaks}{}%
  \urlstyle{tt}}

% Tables: boxed grid as in the proposal, smaller text
\usepackage{etoolbox}
\AtBeginEnvironment{longtable}{\footnotesize\renewcommand{\arraystretch}{1.2}}
\setlength{\LTpre}{6pt}
\setlength{\LTpost}{8pt}
\setlength{\arrayrulewidth}{0.4pt}

% Code blocks: long lines wrap; diagrams stay in one piece
\usepackage{framed}
\usepackage{fvextra}
\RecustomVerbatimEnvironment{verbatim}{Verbatim}{fontsize=\footnotesize, breaklines}
\makeatletter
\@ifundefined{Shaded}{}{\renewenvironment{Shaded}{\begin{snugshade}}{\end{snugshade}}}
\@ifundefined{Highlighting}{}{\RecustomVerbatimEnvironment{Highlighting}{Verbatim}{commandchars=\\\{\}, fontsize=\footnotesize, breaklines, breakanywhere}}
\makeatother

% Figures drawn with TikZ
\usepackage{tikz}
\usetikzlibrary{arrows.meta,positioning,fit,calc,backgrounds,shapes.geometric}
\tikzset{
  ck/.style={font=\footnotesize, >={Stealth[length=5pt]}, line width=0.6pt},
  box/.style={draw=ckline, fill=white, rounded corners=2pt, align=center, inner sep=4pt, minimum height=7mm},
  mod/.style={box, fill=ckgrey},
  this/.style={box, fill=cklight, draw=ckblue, line width=0.9pt},
  group/.style={draw=ckline, dashed, rounded corners=3pt, inner sep=6pt},
  lbl/.style={font=\scriptsize, align=center, inner sep=1.5pt},
  msg/.style={->, draw=ckline},
}
\usepackage{caption}
\captionsetup[figure]{font=small, labelfont=bf, skip=4pt}
\newcommand{\ckfigure}[2]{\begin{figure}[htbp]\centering #1\caption{#2}\end{figure}}

\usepackage{enumitem}
\setlist{itemsep=1pt, topsep=3pt}
\setlength{\parskip}{4pt}
\setlength{\parindent}{0pt}
\raggedbottom
% ---- end of theme ----
\newcommand{\ckmodule}{M6 LED Board and ESP32}
\usepackage{bookmark}
\IfFileExists{xurl.sty}{\usepackage{xurl}}{} % add URL line breaks if available
\urlstyle{same}
% fallback for those not using the hyperref driver hyperxmp:
\makeatletter
\@ifundefined{xmpquote}{\newcommand{\xmpquote}[1]{#1}}{}
\makeatother
\hypersetup{
  pdftitle={M6 LED Board and ESP32},
  pdflang={en},
  colorlinks=true,
  linkcolor={black},
  filecolor={Maroon},
  citecolor={Blue},
  urlcolor={black},
  pdfcreator={LaTeX}}

\title{M6 LED Board and ESP32 --- Module Documentation}
\author{}
\date{}


% ---- colours and drawings ----
\usepackage{colortbl}
\usepackage{float}
\makeatletter
\newcommand{\needspace}[1]{\begingroup\setlength{\dimen@}{#1}\vskip\z@\@plus\dimen@\penalty -100\vskip\z@\@plus -\dimen@\vskip\dimen@\penalty 9999\vskip -\dimen@\vskip\z@skip\endgroup}
\makeatother
% ---- colours ----
\definecolor{srv}{HTML}{2E6FB7}\colorlet{srvL}{srv!15}
\definecolor{cli}{HTML}{2E9E5B}\colorlet{cliL}{cli!15}
\definecolor{adm}{HTML}{E08A1E}\colorlet{admL}{adm!18}
\definecolor{brd}{HTML}{7B4FB3}\colorlet{brdL}{brd!15}
\definecolor{tcpc}{HTML}{1F5FA8}
\definecolor{udpc}{HTML}{D9531E}
\definecolor{obs}{HTML}{D7263D}
\definecolor{tA}{HTML}{00A6D6}
\definecolor{tB}{HTML}{3BB273}
\definecolor{tC}{HTML}{F2C14E}
\definecolor{tD}{HTML}{9B5DE5}

% ---- styles ----
\tikzset{
  nb/.style={rounded corners=3pt, line width=0.9pt, align=center, inner sep=4pt, minimum height=9mm, font=\footnotesize},
  sv/.style={nb, draw=srv, fill=srvL},
  cl/.style={nb, draw=cli, fill=cliL},
  ad/.style={nb, draw=adm, fill=admL},
  bd/.style={nb, draw=brd, fill=brdL},
  hwb/.style={nb, draw=black!50, fill=black!5},
  tcpa/.style={draw=tcpc, line width=1.5pt},
  udpa/.style={draw=udpc, line width=1.5pt, dash pattern=on 5pt off 2.5pt},
  inna/.style={draw=black!55, line width=1pt},
  al/.style={font=\scriptsize, fill=white, inner sep=1.2pt, rounded corners=1pt},
  grp/.style={draw=black!35, dashed, rounded corners=5pt, inner sep=7pt},
  glbl/.style={font=\scriptsize\bfseries, text=black!65},
  hl/.style={line width=2.6pt},
  st/.style={nb, rounded corners=8pt, minimum width=24mm},
  co/.style={draw=black!60, line width=0.5pt},
  cot/.style={font=\scriptsize, align=left, text=black!80},
}

% ---- system map; argument = module to highlight (M1 ... M6, or none) ----
\newcommand{\systemmap}[1]{%
\tikzset{M1/.style={},M2/.style={},M3/.style={},M4/.style={},M5/.style={},M6/.style={},none/.style={}}%
\tikzset{#1/.style={hl}}%
\begin{tikzpicture}[ck]
\node[cl, M4, minimum width=27mm] (m4) at (0,3) {\textbf{M4}\\Player phones};
\node[ad, M5, minimum width=27mm] (m5) at (0,0.4) {\textbf{M5}\\Admin panel};
\node[sv, M2, minimum width=27mm, minimum height=38mm] (m2) at (4.6,1.7) {\textbf{M2}\\Network\\and sync};
\node[sv, M1, minimum width=27mm, minimum height=38mm] (m1) at (8.4,1.7) {\textbf{M1}\\Game server\\and rules};
\node[sv, M3, minimum width=27mm, minimum height=38mm] (m3) at (12.2,1.7) {\textbf{M3}\\Physics,\\traffic\\and AI};
\node[grp, fit=(m2)(m1)(m3)] (pc) {};
\node[glbl, anchor=south west] at (pc.north west) {MAIN COMPUTER};
\node[bd, M6, minimum width=27mm, minimum height=16mm] (m6) at (4.6,-2.9) {\textbf{M6}\\ESP32 program:\\link and LED parts};
\node[hwb, minimum width=25mm] (led) at (8.4,-2.4) {RGB LED\\strips};
\node[hwb, minimum width=25mm] (sm) at (8.4,-3.6) {Sensor and\\4 motors};
\node[grp, fit=(m6)(led)(sm)] (bb) {};
\node[glbl, anchor=north east] at (bb.south east) {LED CITY BOARD};
\draw[tcpa, <->] (m4.east) -- node[al]{WebSocket} (m4.east -| m2.west);
\draw[tcpa, <->] (m5.east) -- node[al]{WebSocket} (m5.east -| m2.west);
\draw[tcpa, <->] ([xshift=-6mm]m6.north) -- node[al]{TCP} ([xshift=-6mm]m2.south);
\draw[udpa, <->] ([xshift=6mm]m6.north) -- node[al]{UDP} ([xshift=6mm]m2.south);
\draw[inna, <->] (m2) -- (m1);
\draw[inna, <->] (m1) -- (m3);
\draw[inna, ->] (m6.east |- led) -- (led);
\draw[inna, <->] (m6.east |- sm) -- (sm);
\ckleg{-1.3,-5.4}{tcpa/TCP or WebSocket: every message arrives in order, udpa/UDP: fast and the newest message wins, inna/inside one program or chip}
\end{tikzpicture}}

% ---- LED board seen from above (6 x 6 roads, 1 unit = 10 cm) ----
\newcommand{\boardbase}{%
\fill[black!88, rounded corners=3pt] (-0.5,-0.5) rectangle (5.5,5.5);
\foreach \i in {0,...,4}{\foreach \j in {0,...,4}{%
  \fill[black!60] (\i+0.13,\j+0.13) rectangle (\i+0.87,\j+0.87);
  \fill[black!22] (\i+0.24,\j+0.24) rectangle (\i+0.52,\j+0.52);
  \fill[black!30] (\i+0.56,\j+0.30) rectangle (\i+0.78,\j+0.62);
  \fill[black!38] (\i+0.28,\j+0.58) rectangle (\i+0.66,\j+0.78);}}
\foreach \i in {0,...,5}{%
  \draw[black, line width=4.5pt] (\i,-0.05) -- (\i,5.05);
  \draw[black, line width=4.5pt] (-0.05,\i) -- (5.05,\i);}
\foreach \i in {0,...,5}{\foreach \k in {0,...,30}{%
  \fill[black!55] (\i,\k/6) circle (0.025);
  \fill[black!55] (\k/6,\i) circle (0.025);}}
\foreach \x/\y/\c in {0/4/tA, 4/4/tB, 0/0/tD, 4/0/tC}{%
  \fill[\c!30!black!60] (\x+0.13,\y+0.13) rectangle (\x+0.87,\y+0.87);
  \fill[\c] (\x+0.3,\y+0.3) rectangle (\x+0.7,\y+0.7);
  \fill[white] (\x+0.4,\y+0.5) circle (0.05); \fill[white] (\x+0.6,\y+0.5) circle (0.05);}}
\newcommand{\nmark}[3]{\node[circle, fill=yellow!85!orange, draw=black, line width=0.4pt, inner sep=0.6pt, minimum size=3.6mm, font=\tiny\bfseries] at (#1,#2) {#3};}
\newcommand{\bike}[4]{%
\fill[#3!35] ($(#1,#2)+(#4:0.40)$) circle (0.05);
\fill[#3!65] ($(#1,#2)+(#4:0.21)$) circle (0.075);
\fill[#3, draw=white, line width=0.4pt] (#1,#2) circle (0.12);}
\newcommand{\dpoint}[3]{\draw[#3!50, line width=0.8pt] (#1,#2) circle (0.26); \draw[#3, line width=1.4pt] (#1,#2) circle (0.16); \fill[#3] (#1,#2) circle (0.07);}
\newcommand{\heavy}[2]{\draw[obs!60!black, line width=4pt] (#1) -- (#2); \draw[obs, line width=4pt, dash pattern=on 3pt off 5pt] (#1) -- (#2);}
\newcommand{\roadwork}[2]{\draw[obs, line width=4pt] (#1) -- (#2);}
\newcommand{\accident}[2]{\draw[obs, line width=4pt, dash pattern=on 3pt off 3pt] (#1) -- (#2);}
\newcommand{\puddle}[2]{\draw[cyan!55, line width=4pt] (#1) -- (#2); \draw[cyan!15, line width=1.6pt, dash pattern=on 1pt off 2.5pt] (#1) -- (#2);}
\newcommand{\findme}[2]{\fill[white, opacity=0.25] (#1,#2) circle (0.36); \fill[white, opacity=0.5] (#1,#2) circle (0.24); \fill[white] (#1,#2) circle (0.13);}

% ---- phone game screen (7 x 3.4 units) ----
\newcommand{\phonemock}{%
\fill[black!85, rounded corners=9pt] (-0.3,-0.3) rectangle (7.3,3.7);
\fill[cyan!20] (0,1.7) rectangle (7,3.4);
\fill[cli!14] (0,0) rectangle (7,1.7);
\filldraw[fill=srv!30, draw=srv!70!black] (0.2,1.5) rectangle (1.0,3.0);
\filldraw[fill=adm!35, draw=adm!70!black] (1.1,1.6) rectangle (1.8,2.4);
\filldraw[fill=brd!30, draw=brd!70!black] (1.9,1.65) -- (2.6,1.65) -- (2.6,2.2) -- (2.25,2.5) -- (1.9,2.2) -- cycle;
\filldraw[fill=adm!35, draw=adm!70!black] (4.5,1.65) rectangle (5.1,2.3);
\filldraw[fill=srv!30, draw=srv!70!black] (5.2,1.55) rectangle (6.0,3.1);
\filldraw[fill=cli!35, draw=cli!70!black] (6.1,1.6) rectangle (6.85,2.6);
\fill[brown!60!black] (4.17,1.6) rectangle (4.23,1.8); \fill[cli!70!black] (4.2,1.95) circle (0.18);
\fill[black!55] (1.8,0) -- (5.2,0) -- (3.75,1.7) -- (3.25,1.7) -- cycle;
\draw[white, line width=0.8pt, dash pattern=on 4pt off 4pt] (3.5,0) -- (3.5,1.7);
\filldraw[fill=black!15, draw=black!70] (3.62,1.12) rectangle (3.88,1.32);
\fill[obs] (3.05,1.22) -- (3.15,1.48) -- (3.25,1.22) -- cycle;
\fill[black!85] (3.43,0.08) rectangle (3.57,0.3);
\fill[tA] (3.3,0.28) rectangle (3.7,0.6);
\fill[brown!65] (3.34,0.6) rectangle (3.66,0.82);
\fill[tA!60!black] (3.5,0.95) circle (0.13);
\node[fill=black!70, text=white, rounded corners=2pt, font=\scriptsize\bfseries, inner sep=2pt] at (3.5,3.15) {2:41};
\foreach \c/\s/\x in {tA/120/0.35,tB/95/0.95,tC/140/1.55,tD/80/2.15}
  \node[fill=\c, text=white, font=\tiny\bfseries, rounded corners=1pt, inner sep=1.5pt, minimum width=5mm] at (\x,3.15) {\s};
\fill[white, opacity=0.9] (6.45,2.85) circle (0.38); \draw[black!60] (6.45,2.85) circle (0.38);
\draw[-{Stealth[length=5pt]}, line width=1.8pt, obs] (6.45,2.85) ++(230:0.2) -- ++(50:0.5);
\fill[white, opacity=0.85] (0.55,0.5) circle (0.32); \fill[black!70] (0.43,0.5) -- (0.65,0.63) -- (0.65,0.37) -- cycle;
\fill[white, opacity=0.85] (1.4,0.5) circle (0.32); \fill[black!70] (1.52,0.5) -- (1.3,0.63) -- (1.3,0.37) -- cycle;
\node[fill=obs!80, text=white, rounded corners=3pt, font=\tiny\bfseries, minimum width=8mm, minimum height=6mm] at (5.6,0.5) {Brake};
\node[fill=cli, text=white, rounded corners=3pt, font=\tiny\bfseries, minimum width=8mm, minimum height=8mm] at (6.5,0.6) {Gas};
\node[fill=white, text=black, rounded corners=3pt, font=\tiny\bfseries, minimum width=9mm, minimum height=5mm] at (0.9,1.45) {Find me};}

% ---- side cut of the board (14 units wide) ----
\newcommand{\boardside}{%
\fill[brown!20] (-0.6,-0.35) rectangle (14.6,0);
\foreach \x/\h/\c in {1.0/0.7/black!18, 2.2/1.1/black!22, 4.6/1.9/black!25, 5.8/2.6/black!28, 7.4/2.2/black!25, 8.6/1.6/black!22, 10.9/1.0/black!18}
  \filldraw[fill=\c, draw=black!45] (\x,1.5) rectangle (\x+0.9,1.5+\h);
\filldraw[fill=tB!35, draw=tB!70!black] (12.4,1.5) rectangle (13.4,2.3);
\fill[black!70, rounded corners=1pt] (0.5,0) rectangle (1.1,0.9);
\fill[black!70, rounded corners=1pt] (12.9,0) rectangle (13.5,0.9);
\filldraw[fill=brown!45, draw=brown!70!black] (0,0.9) rectangle (14,1.3);
\fill[black!80] (0,1.3) rectangle (14,1.38);
\foreach \x in {0.3,0.9,...,13.8} \fill[black!35] (\x,1.42) circle (0.06);
\foreach \x/\c in {3.3/tA, 3.9/tA, 6.9/obs, 7.5/obs, 10.5/tC} \fill[\c] (\x,1.42) circle (0.07);
\fill[white, opacity=0.55] (0,1.38) rectangle (14,1.52);
\filldraw[fill=brd!40, draw=brd] (2.3,0.35) rectangle (4.3,0.85);
\node[font=\tiny] at (3.3,0.6) {ESP32};
\filldraw[fill=brd, draw=brd!60!black] (6.75,0.62) rectangle (7.25,0.88);
\filldraw[fill=adm!60, draw=adm!70!black] (1.8,0.6) circle (0.24);
\filldraw[fill=adm!60, draw=adm!70!black] (12.2,0.6) circle (0.24);
\filldraw[fill=black!15, draw=black!50] (9.3,0.35) rectangle (10.9,0.85);
\node[font=\tiny] at (10.1,0.6) {Power};}

% ---- layout helpers ----
\newenvironment{ckfig}{\begin{figure}[!htb]\centering}{\end{figure}}
\renewcommand{\topfraction}{0.9}\renewcommand{\bottomfraction}{0.8}\renewcommand{\textfraction}{0.07}\renewcommand{\floatpagefraction}{0.85}
\setcounter{topnumber}{3}\setcounter{bottomnumber}{2}\setcounter{totalnumber}{4}
\newcommand{\msub}[1]{\par\needspace{5\baselineskip}\smallskip\noindent{\bfseries #1:}\par\nobreak\vspace{-2pt}}
\tikzset{msga/.style={draw=cli!70!black, line width=0.8pt}, cmda/.style={draw=obs, line width=1pt}, loopa/.style={draw=srv, line width=1pt}, hwa/.style={draw=black!60, line width=1.2pt, dash pattern=on 1.5pt off 1.5pt}, wanta/.style={draw=cli, line width=1.6pt}, reala/.style={draw=obs, line width=1.6pt}, siga/.style={draw=srv, line width=0.8pt}, thra/.style={draw=obs, line width=0.9pt, dash pattern=on 4pt off 2pt}, evta/.style={draw=cli, line width=1pt}}
\tikzset{iol/.style={font=\scriptsize, fill=white, inner sep=1.2pt, rounded corners=1pt, align=center}}
% legend row: start point, list of style/text
\newcommand{\ckleg}[2]{%
\coordinate (lgp) at (#1);
\foreach \ls/\lt in {#2}{%
  \draw[\ls,->] (lgp) -- ++(0.8,0) node[right, font=\scriptsize, text=black] (lgt) {\lt};
  \coordinate (lgp) at ([xshift=5mm]lgt.east);}}
% purpose table: colour (unused), runs on, technology, requirements
\newcommand{\ckinfo}[4]{\par\smallskip\noindent{\renewcommand{\arraystretch}{1.25}%
\begin{tabular}{|L{0.30}|L{0.50}|L{0.20}|}\hline
\textbf{Runs on} & \textbf{Technology} & \textbf{Requirements} \\ \hline
#2 & #3 & #4 \\ \hline
\end{tabular}}\par\smallskip}
% module data map: n in, n out, centre text, centre style, inputs, outputs, legend list
% inputs and outputs: {name}/style/{data}/arrow style
\newcommand{\iomap}[7]{%
\begin{tikzpicture}[ck]
\pgfmathsetmacro{\nmax}{max(#1,#2)}
\node[#4, line width=2pt, minimum width=30mm, minimum height=\nmax*10.5mm+2mm, font=\normalsize] (c) at (0,0) {#3};
\foreach \nm/\st/\dd/\ar [count=\ii] in {#5}{%
  \pgfmathsetmacro{\yy}{((#1+1)/2-\ii)*1.05}
  \node[\st, minimum width=25mm, minimum height=8.5mm, font=\scriptsize] (in\ii) at (-6.6,\yy) {\nm};
  \draw[\ar,->] (in\ii.east) -- node[iol]{\dd} (c.west |- in\ii);}
\foreach \nm/\st/\dd/\ar [count=\ii] in {#6}{%
  \pgfmathsetmacro{\yy}{((#2+1)/2-\ii)*1.05}
  \node[\st, minimum width=25mm, minimum height=8.5mm, font=\scriptsize] (out\ii) at (6.6,\yy) {\nm};
  \draw[\ar,->] (c.east |- out\ii) -- node[iol]{\dd} (out\ii.west);}
\node[font=\scriptsize\bfseries, text=black!60] at (-6.6,\nmax*0.525+0.4) {Inputs come from};
\node[font=\scriptsize\bfseries, text=black!60] at (6.6,\nmax*0.525+0.4) {Outputs go to};
\ckleg{-7.6,-\nmax*0.525-0.5}{#7}
\end{tikzpicture}}

% keep an intro sentence and its table together on one page
\newcommand{\ckkeep}[1]{\par\noindent\begin{minipage}{\linewidth}\setlength{\parskip}{4pt}#1\end{minipage}\par\smallskip}
\widowpenalty=10000 \clubpenalty=10000

\begin{document}

\thispagestyle{ckfirst}
\begin{center}
{\bfseries GEBZE TECHNICAL UNIVERSITY\par}
{Computer Engineering Department\par}
{CSE 396 Module Documentation\par}
\vspace{8pt}
{\LARGE\bfseries Çılgın Kurye\par}
\vspace{3pt}
{\large\itshape Module M6: LED Board and ESP32\par}
\vspace{6pt}
\textbf{Company:} Board\&Byte\hspace{1.4em}\textbf{Group:} Group-8\hspace{1.4em}\textbf{Date:} 7 November 2026
\end{center}
\begin{center}\footnotesize\renewcommand{\arraystretch}{1.15}
\begin{tabular}{|p{0.27\linewidth}|p{0.15\linewidth}|p{0.27\linewidth}|p{0.15\linewidth}|}\hline
\textbf{Name} & \textbf{Student No.} & \textbf{Name} & \textbf{Student No.} \\ \hline
Tuğberk Akbulut & 230104004070 & Hasan Alperen Elmas & 1801042604 \\ \hline
Hasan Hüseyin Özdemir & 230104004048 & Muhib Taha Boy & 230104004165 \\ \hline
Ayşe Poyraz & 230104004085 & Umutcan Ocak & 220104004903 \\ \hline
\end{tabular}
\end{center}
\vspace{2pt}



% ---- table and list helpers ----
\newcolumntype{L}[1]{>{\raggedright\arraybackslash}p{\dimexpr#1\linewidth-2\tabcolsep-\arrayrulewidth\relax}}
\newcolumntype{C}[1]{>{\centering\arraybackslash}p{\dimexpr#1\linewidth-2\tabcolsep-\arrayrulewidth\relax}}
\newcommand{\ckdone}{\item[$\boxtimes$]}
\newcommand{\ckopen}{\item[$\square$]}
\newcommand{\mdot}{$\bullet$}
\newcounter{msec}
\newcounter{chap}
% ---- end of helpers ----
\newcommand{\msection}[1]{\needspace{8\baselineskip}\stepcounter{msec}\section{\arabic{msec}. #1}}

\msection{Purpose}

M6 is the physical city board and the program on its ESP32. The board is 60 × 60 cm. Programmable
addressable RGB LED strips lie in all roads and in the four stores, and there are 3D printed buildings,
a motion sensor and four vibration motors. In the game the board is the only map, so it must be clear
to read from every side of the table. One ESP32 runs one Arduino program with two parts. The link part
gets the game state from the main computer over Wi-Fi, finds the strength and the place of a hit on the
board, drives the vibration motors and reports the health of the board to the admin. The LED part turns
the game state into light.

\ckinfo{brd}{ESP32 in the LED board, and the board itself for the hardware}{C++ with Arduino for ESP32, TCP and UDP over Wi-Fi with MessagePack messages, programmable addressable RGB LED strips with an external power supply, a motion sensor with accelerometer and gyroscope, 4 vibration motors, 3D printing, a plywood base}{R2 R5 R6 R7 R8 R10 R11}

M6 talks to the server over Wi-Fi and to the hardware under the board. The figure shows the data on
both sides.

\begin{ckfig}
\iomap{5}{4}{\textbf{M6}\\LED board\\and ESP32}{bd}
{{M2, TCP}/sv/{config, vibrate}/tcpa, {M2, UDP}/sv/{board state,\\delay probe}/udpa, {Motion sensor}/hwb/{sensor data}/hwa, {USB cable}/hwb/{Wi-Fi and\\calibration settings}/hwa, {City plan file}/hwb/{road grid,\\store places}/inna}
{{M2, TCP}/sv/{earthquake event,\\status}/tcpa, {M2, UDP}/sv/{sensor samples,\\delay replies}/udpa, {4 motors}/hwb/{vibration power}/hwa, {RGB LED strips}/hwb/{colour of every LED,\\60 frames per second}/hwa}
{tcpa/TCP, udpa/UDP, inna/built into the program, hwa/wire to the hardware}
\end{ckfig}

\msection{Requirements covered}

\ckkeep{%
M6 takes part in seven requirements. The last column says what will work when M6 is finished.

\begin{tabular}{|L{0.06}|L{0.15}|L{0.40}|L{0.39}|}
\hline \textbf{ID} & \textbf{Requirement} & \textbf{What M6 does for it} & \textbf{When M6 is done} \\ \hline 
R2 & Delivery cycle & Delivery points blink in the player's colour, a flash when the package is delivered & Each player finds the own delivery point on the board \\ \hline
R5 & No map on the phone & The board is the map with all motorcycles. ``Find me'' makes the player's dot white for 1 s & A player can always find the own motorcycle \\ \hline
R6 & Obstacles & A red pattern for each type on the whole road segment, light blue for a puddle & Every obstacle type is clear from a distance \\ \hline
R7 & Earthquake & Finds hand hits and motor shaking with the sensor, with a strength from 0 to 10 and a zone in a 3 × 3 grid. Drives the motors for the admin button. A shock wave from the earthquake centre & A hit on the table starts an earthquake in the game, and the earthquake is seen on the table \\ \hline
R8 & Physical LED board & 60 × 60 cm, 6 roads in each direction, LEDs on all roads, printed buildings, four stores. One ESP32 for the LEDs, the sensor and the motors brings the game state over Wi-Fi & The finished board stands on the table and shows the live game without a cable to the computer \\ \hline
R10 & Network and measurement & Delay replies, sensor samples for the seismic graph, comes back after a lost link & The admin sees the board delay and the sensor signal live \\ \hline
R11 & Engineering requirements & PC emulators for both sides of the link, a demo mode without a server & The server and the board can be built and tested without each other \\ \hline
\end{tabular}}

\msection{Inputs and outputs}

\ckkeep{%
Messages that must arrive use TCP, fast and repeated data uses UDP. The board state comes 10 to 30 times
per second, but the LED part draws 60 frames per second, so it fills the time between two states by
itself.

\begin{tabular}{|L{0.11}|L{0.17}|L{0.18}|L{0.54}|}
\hline \textbf{Direction} & \textbf{Data} & \textbf{From / to} & \textbf{Content} \\ \hline 
In & Board state & M2, UDP & 10 to 30 per second: phase, players (team, colour, position, flags), delivery points, obstacles, earthquake \\ \hline
In & Config & M2, TCP & Player colours, brightness from 0.1 to 1, power limit \\ \hline
In & Vibrate & M2, TCP & Motors to use, strength from 0 to 1, duration \\ \hline
In & Delay probe & M2, UDP & A small message with the server time, answered at once \\ \hline
In & Sensor & Hardware & Motion data from the accelerometer and the gyroscope \\ \hline
In & Settings & USB cable & Wi-Fi name and password, backup server address, calibration values \\ \hline
In & City plan & Shared file & Road grid and store places, turned into the LED map when the code is built \\ \hline
Out & Light & LED strips & The colour of every LED, 60 frames per second \\ \hline
Out & Vibration & 4 motors & Motor power with a soft start and safety limits \\ \hline
Out & Earthquake event & M2, TCP, to M1 and M5 & Source (hand or motors), strength from 0 to 10, centre, zone from 0 to 8, confidence \\ \hline
Out & Sensor samples & M2, UDP, to M5 & Sensor data for the seismic graph, sent in small groups \\ \hline
Out & Status, delay reply & M2 & Status on TCP once per second: frame rate, Wi-Fi signal, power estimate, errors. Delay replies on UDP \\ \hline
\end{tabular}}

\msection{Main components}

\ckkeep{%
M6 has eleven parts. The first six are the link part and the last five are the LED part. All of them
run in the main loop of the Arduino program.

\begin{tabular}{|L{0.21}|L{0.39}|L{0.40}|}
\hline \textbf{Component} & \textbf{Takes} & \textbf{Gives} \\ \hline 
Network link & Wi-Fi, server address & One TCP channel and one UDP stream, a new link after a loss \\ \hline
Message codec & Raw messages & Checked messages. Old, broken or wrong version messages are dropped \\ \hline
State hand over & Checked board state and config & The newest full state for the LED part \\ \hline
Hit detector & Sensor samples & Earthquake events with strength, zone and confidence \\ \hline
Motor controller & Vibrate commands & Motor power with a soft start and safety limits \\ \hline
Status and settings & Frame report, USB settings & Status reports, delay replies, settings saved in flash \\ \hline
LED map & City plan & The place of every LED on the roads and in the stores, made by a script \\ \hline
Motion smoother & The last two board states & Smooth motorcycle positions along the roads for every frame \\ \hline
Layers and effects & Board state, phase, link state & Five layers and the patterns for obstacles, earthquakes and game phases \\ \hline
Post processing and output & The mixed frame, config & Corrected colours inside the brightness and power limit, sent to the strips \\ \hline
Demo mode & A built in scenario & A full game look without a server \\ \hline
\end{tabular}}

\msection{How it works}

\msub{One program with a main loop}
The ESP32 runs one Arduino program. The main loop repeats the same steps again and again: it reads the
new network messages, reads the sensor, drives the motors, lets the LED part draw one LED frame and
sends the reports. The link part keeps the newest full board state in memory. At the start of each
frame the LED part reads this state, so it always draws the latest game.

\begin{ckfig}
\begin{tikzpicture}[ck]
\draw[brd, line width=1.2pt, rounded corners=6pt, fill=brd!4] (2.6,-1.2) rectangle (16.0,1.55);
\node[font=\scriptsize\bfseries, text=brd, anchor=north west] at (2.7,1.5) {ESP32: one Arduino program};
\node[bd, minimum width=19mm] (n) at (4.1,0.2) {Read network\\messages};
\node[bd, minimum width=19mm] (s) at (6.75,0.2) {Read sensor,\\find hits};
\node[bd, minimum width=19mm] (m) at (9.4,0.2) {Drive\\motors};
\node[nb, draw=brd, fill=brd!25, line width=1.3pt, minimum width=19mm] (l) at (12.05,0.2) {Draw\\LED frame};
\node[bd, minimum width=19mm] (r) at (14.7,0.2) {Send\\reports};
\foreach \u/\v in {n/s,s/m,m/l,l/r} \draw[inna,->] (\u) -- (\v);
\draw[loopa,->] (r.south) -- ++(0,-0.55) -| node[al, pos=0.25]{next round of the loop} (n.south);
\node[sv, minimum width=20mm, minimum height=22mm] (pc) at (0,0.2) {Main\\computer\\(M2)};
\draw[tcpa,<->] ([yshift=5mm]pc.east) -- node[iol]{TCP} ([yshift=5mm]pc.east -| n.west);
\draw[udpa,<->] ([yshift=-5mm]pc.east) -- node[iol]{UDP} ([yshift=-5mm]pc.east -| n.west);
\node[hwb, minimum width=19mm] (hs) at (6.75,-2.4) {Motion sensor};
\node[hwb, minimum width=19mm] (hm) at (9.4,-2.4) {4 motors};
\node[hwb, minimum width=19mm] (hl) at (12.05,-2.4) {RGB LED strips};
\draw[hwa,->] (hs.north) -- (hs.north |- 0,-1.2);
\draw[hwa,->] (hm.north |- 0,-1.2) -- (hm.north);
\draw[hwa,->] (hl.north |- 0,-1.2) -- (hl.north);
\ckleg{0,-3.3}{tcpa/TCP, udpa/UDP, inna/steps of the loop, loopa/next round, hwa/wire to the hardware}
\end{tikzpicture}
\end{ckfig}

\msub{The board}
The base is plywood. On it there are 6 roads in each direction, so the city has 5 × 5 blocks. The LED
strips lie in the roads between the crossings. There are no LEDs at the crossings, because two strips
cannot lie on top of each other. The buildings are 3D printed from the same models as the 3D city on
the phone, so a player finds the same building on the board and on the phone. The LEDs need an
external power supply, and the software keeps the power under the limit in the config.

\begin{ckfig}
\begin{tikzpicture}[x=1.15cm, y=1.15cm]
\boardbase
\bike{1}{3.5}{tA}{90}\bike{2.5}{4}{tA}{180}
\bike{3.5}{4}{tB}{0}\bike{5}{2.5}{tB}{90}
\bike{1.5}{1}{tD}{180}\bike{0}{1.5}{tD}{-90}
\bike{4}{2.5}{tC}{-90}\bike{4.5}{0}{tC}{180}
\dpoint{2.5}{2}{tA}\dpoint{5}{3.5}{tB}\dpoint{0.5}{3}{tD}\dpoint{3}{0.5}{tC}
\heavy{0,2}{1,2}
\roadwork{3,2}{3,3}
\accident{4,1}{5,1}
\puddle{2,4}{2,5}
\findme{3.5}{4}
\nmark{0.5}{5.25}{1}\nmark{0.62}{3.5}{2}\nmark{2.5}{2.45}{3}\nmark{0.5}{2.3}{4}\nmark{3.35}{2.5}{5}\nmark{4.5}{1.35}{6}\nmark{1.65}{4.5}{7}\nmark{3.5}{4.5}{8}
\node[anchor=west, font=\small, align=left] at (6.0,2.5) {\begin{tabular}{@{}c@{\ \ }l@{}}
\textbf{1} & Team store in the team colour\\[3pt]
\textbf{2} & Motorcycle: a bright dot with a short tail\\[3pt]
\textbf{3} & Delivery point: blinks in the player's tone\\[3pt]
\textbf{4} & Heavy traffic: a slow dim red wave\\[3pt]
\textbf{5} & Road work: bright steady red\\[3pt]
\textbf{6} & Accident: red blinking\\[3pt]
\textbf{7} & Water puddle: light blue shimmer\\[3pt]
\textbf{8} & ``Find me'': white flash for 1 second\end{tabular}};
\end{tikzpicture}
\end{ckfig}

\msub{What is inside the board}
The figure below is a side cut of the board. The motion sensor is fixed at the centre, under the
plywood, so it feels a hit from every side. The four vibration motors are under the four corners. The
ESP32 and a power supply for the LEDs and the motors sit under the board too.

\begin{ckfig}
\begin{tikzpicture}[x=0.9cm, y=0.9cm]
\boardside
\nmark{3.3}{-0.15}{1}\nmark{7.0}{0.3}{2}\nmark{1.8}{0.05}{3}\nmark{12.2}{0.05}{3}\nmark{10.1}{-0.15}{4}
\nmark{6.25}{1.7}{5}\nmark{6.25}{4.4}{6}\nmark{12.9}{2.55}{7}
\node[anchor=north, font=\small] at (7,-0.6) {\begin{tabular}{@{}c@{\ \ }l@{\qquad}c@{\ \ }l@{\qquad}c@{\ \ }l@{}}
\textbf{1} & ESP32 & \textbf{4} & Power supply & \textbf{6} & 3D printed buildings \\
\textbf{2} & Motion sensor at the centre & \textbf{5} & RGB LED strips on the roads & \textbf{7} & Team store with LEDs \\
\textbf{3} & Vibration motors at the corners & & & & \end{tabular}};
\end{tikzpicture}
\end{ckfig}

\msub{Colours}
There are four teams and each team has one colour: blue, green, yellow and purple. The two players of a
team have a light and a dark tone of the team colour. Red is used only for obstacles, so red always
means ``this road slows you down''. The pattern of the red tells the type of the obstacle. A water
puddle is light blue, because it makes the road slippery, not slow. White is used only for ``Find
me'', so a player always knows that the white dot is the own motorcycle.

\begin{ckfig}
\begin{tikzpicture}[ck]
\foreach \r/\t in {0/Heavy traffic, 1/Road work, 2/Accident, 3/Water puddle}
  \node[font=\scriptsize\bfseries, anchor=east] at (2.3,-\r*0.75) {\t};
\foreach \k in {0,...,5}{%
  \pgfmathsetmacro{\a}{40+40*sin(\k*60)}
  \fill[obs!\a!black] (2.7+\k*0.32,0) circle (0.11);
  \fill[obs] (2.7+\k*0.32,-0.75) circle (0.11);
  \pgfmathtruncatemacro{\w}{mod(\k,2)}
  \ifnum\w=0 \fill[obs] (2.7+\k*0.32,-1.5) circle (0.11); \else \fill[black!75] (2.7+\k*0.32,-1.5) circle (0.11); \fi
  \pgfmathsetmacro{\b}{30+20*sin(\k*97)}
  \pgfmathsetmacro{\b}{45+25*sin(\k*97)}\fill[cyan!\b!white] (2.7+\k*0.32,-2.25) circle (0.11);}
\node[font=\scriptsize, anchor=west] at (4.75,0) {the light moves slowly along the road};
\node[font=\scriptsize, anchor=west] at (4.75,-0.75) {the whole road is red, it does not move};
\node[font=\scriptsize, anchor=west] at (4.75,-1.5) {red blinks on and off};
\node[font=\scriptsize, anchor=west] at (4.75,-2.25) {light blue light that shimmers};
\foreach \i/\c in {0/tA,1/tB,2/tC,3/tD}{%
  \fill[\c] (11.0+\i*0.75,-0.55) circle (0.17); \fill[\c!55] (11.0+\i*0.75,-1.05) circle (0.17);}
\node[font=\scriptsize\bfseries] at (12.12,0.0) {Team colours};
\node[font=\scriptsize] at (12.12,-1.5) {light and dark tone};
\end{tikzpicture}
\end{ckfig}

\msub{Layers}
The LED part draws 60 frames per second. Each frame is made from five layers, from the bottom to the
top: the roads, the obstacles, the delivery points, the motorcycles and a top layer for short effects.
A higher layer covers a lower one, so when two things fall on the same LED, the higher one is shown.
When two motorcycles are on the same LED, their colours take turns. After the layers, the LED part
corrects the colours for the eye, sets the brightness from the config, and lowers all LEDs a little if
the frame would need more power than the limit.

\begin{ckfig}
\begin{tikzpicture}[ck]
\foreach \i/\t/\c in {0/Roads/gray, 1/Obstacles/obs, 2/Delivery points/adm, 3/Motorcycles/tA, 4/Top effects/brd}{%
  \fill[\c!25, draw=\c, line width=0.9pt] (\i*0.35,\i*0.62) -- ++(4.2,0) -- ++(0.9,0.45) -- ++(-4.2,0) -- cycle;
  \node[font=\scriptsize, anchor=west] at (\i*0.35+5.3,\i*0.62+0.22) {\t};}
\draw[inna,->] (8.5,1.9) -- node[al]{mix} (9.6,1.9);
\node[nb, draw=srv, fill=srvL, minimum width=20mm] (pp) at (11.0,1.9) {Colour,\\brightness,\\power limit};
\draw[hwa,->] (pp) -- (12.8,1.9);
\ckleg{0,-0.6}{inna/inside the program, hwa/data wire to the strips}
\node[nb, draw=black!50, fill=black!8, minimum width=14mm] at (13.6,1.9) {RGB LED\\strips};
\end{tikzpicture}
\end{ckfig}

\msub{Smooth motion on the roads}
The board state comes less often than the frames, so the motion smoother moves each motorcycle dot
between the last two states. It moves the dot along the roads, through the crossing, and never in a
straight line over a block. The dot is drawn one state behind the newest one, so it adds a small
delay. If a state is lost, the dot keeps moving along its road with its last speed until the next
state comes. The dot has a short tail, so the direction of motion is easy to see.

\msub{Two channels}
The board keeps one TCP channel and one UDP stream to the server. Messages that must arrive go on TCP:
the config, the vibrate commands, the earthquake events and the status report. Each TCP message starts
with its length, because TCP is only a stream of bytes. Data that comes again and again goes on UDP:
the board state, the sensor samples and the delay probes. If one UDP packet is lost, the next one
replaces it, so the LEDs never wait. Every message has an order number, so the board drops a UDP
packet that comes late. Each start of the server makes a new session number. When the board sees a
new session number, it starts its order numbers again, so it does not drop new messages as old ones.

\msub{Finding a hit}
The sensor measures the motion all the time. First a filter takes away gravity and slow
tilting, so only the shaking is left. The detector also learns the normal noise of the room and puts the
start threshold well above it. So the threshold goes up on a noisy table, and normal play does not
start an earthquake. When the signal passes the threshold, a hit starts, and the detector takes a
short window from this moment. The stronger the shaking in this window, the higher the strength. Like
a real earthquake scale, the strength grows slowly for very strong hits, so a light tap and a hard
hit both get a useful value. The scale is set once in calibration with a standard hit, and the result
is kept between 0 and 10.

\msub{Finding the zone}
The accelerometer shows the direction of the first push of the board, and the gyroscope shows how the
board tilts. Together they show the side of the hit, and the size of the push shows how far the hit is
from the centre. From this the detector picks one of 9 zones in a 3 × 3 grid and gives a confidence
value. After a hit the detector waits a short time, so one hit makes only one event.

\begin{ckfig}
\begin{tikzpicture}[ck]
\begin{scope}
\fill[adm!15] (3.0,-1.7) rectangle (4.5,1.9);
\draw[black!40] (-0.2,0) -- (8.2,0);
\draw[black!40,->] (0,-1.7) -- (0,2.0);
\draw[obs, line width=0.9pt, dash pattern=on 4pt off 2pt] (0,0.45) -- (8,0.45);
\draw[obs, line width=0.9pt, dash pattern=on 4pt off 2pt] (0,-0.45) -- (8,-0.45);
\draw[srv, line width=0.8pt, smooth, samples=400, domain=0:8] plot (\x, {0.12*sin(1500*\x) + 0.06*sin(1900*\x) + (\x>3 ? 1.6*exp(-1.6*max(\x-3,0))*sin(2200*(\x-3)) : 0)});
\node[font=\scriptsize, text=adm!80!black] at (3.75,2.15) {short window};
\node[font=\scriptsize, text=obs, anchor=west] at (5.6,0.7) {threshold};
\node[font=\scriptsize, anchor=west] at (0.3,-1.1) {room noise};
\node[font=\scriptsize] at (4.2,-2.0) {time};
\node[al, draw=cli, text=cli!70!black, line width=0.8pt] at (6.9,1.5) {event sent};
\draw[evta, ->] (4.5,1.5) -- (6.25,1.5);
\end{scope}
\begin{scope}[shift={(9.4,-1.5)}]
\foreach \i in {0,1,2}{\foreach \j in {0,1,2}{%
  \pgfmathtruncatemacro{\z}{(2-\j)*3+\i}
  \filldraw[fill=black!6, draw=black!40] (\i,\j) rectangle ++(1,1);
  \node[font=\tiny, text=black!50] at (\i+0.18,\j+0.82) {\z};}}
\fill[obs!45] (2,1) rectangle ++(1,1);
\node[font=\tiny, text=black!70] at (2.18,1.82) {5};
\fill[brd] (1.5,1.5) circle (0.09);
\draw[obs, line width=1pt, ->] (1.5,1.5) -- (2.5,1.45);
\fill[obs] (2.55,1.45) circle (0.1);
\node[font=\scriptsize] at (1.5,-0.3) {zone 5, confidence 0.8};
\node[font=\scriptsize] at (1.5,3.3) {board from above};
\end{scope}
\ckleg{0,-2.6}{siga/sensor signal, thra/start threshold, evta/event to the server}
\end{tikzpicture}
\end{ckfig}

\msub{Motors}
The four motors are under the corners, at zones 0, 2, 6 and 8. When the admin presses the earthquake
button, M1 sends a vibrate command on TCP. The command says which motors to use, the strength and the
time. When the admin chooses one zone, M1 uses the motors nearest to it. For the whole board it uses all
four. M6 starts the motors slowly, so the power does not jump. While the motors run, the detector stops
learning the room noise, and the whole shaking gives only one event with ``motors'' as its source. So
the earthquake in the game follows the real shaking of the table, and the motors do not start a second
earthquake.

\msub{Earthquake on the board}
When an earthquake starts, a ring of light grows from the centre of the earthquake over the whole
board, like a wave on water, and slowly fades. The road work after an earthquake is not a decision of
M6. M1 puts it into the obstacle list, and M6 only draws it like any other obstacle until M1 removes it
from the list.

\begin{ckfig}
\begin{tikzpicture}[x=0.62cm, y=0.62cm]
\foreach \s/\r in {0/0.6, 7.2/1.6, 14.4/2.8}{%
  \begin{scope}[shift={(\s,0)}]
  \boardbase
  \begin{scope}
  \clip (-0.5,-0.5) rectangle (5.5,5.5);
  \draw[white, line width=5pt, opacity=0.35] (3.4,2.6) circle (\r+0.25);
  \draw[yellow!70!orange, line width=3pt] (3.4,2.6) circle (\r);
  \end{scope}
  \fill[obs] (3.4,2.6) circle (0.12);
  \end{scope}}
\node[font=\scriptsize] at (2.5,-1.1) {start};
\node[font=\scriptsize] at (9.7,-1.1) {a moment later};
\node[font=\scriptsize] at (16.9,-1.1) {later};
\end{tikzpicture}
\end{ckfig}

\msub{Game phases and the link}
Every game phase has its own look, so the people around the table always know what happens.

\begin{ckfig}
\begin{tikzpicture}[x=0.4cm, y=0.4cm]
\foreach \sx/\t in {0/Lobby, 7.6/Countdown, 15.2/Running, 22.8/Paused, 30.4/Finished}{%
  \begin{scope}[shift={(\sx,0)}]\boardbase\end{scope}
  \node[font=\scriptsize\bfseries] at (\sx+2.5,-1.3) {\t};}
\begin{scope}[shift={(7.6,0)}]
  \draw[adm!70!yellow, line width=2.2pt] (0,0) rectangle (5,5);
  \draw[adm!50!yellow, line width=2.2pt, opacity=0.6] (1,1) rectangle (4,4);
  \node[font=\bfseries, text=white] at (2.5,2.5) {3};
\end{scope}
\begin{scope}[shift={(15.2,0)}]
  \bike{1}{3.5}{tA}{90}\bike{3.5}{4}{tB}{0}\bike{4}{2.5}{tC}{-90}\bike{1.5}{1}{tD}{180}
  \dpoint{2.5}{2}{tA}\roadwork{3,2}{3,3}
\end{scope}
\begin{scope}[shift={(22.8,0)}]
  \bike{1}{3.5}{tA}{90}\bike{3.5}{4}{tB}{0}\bike{4}{2.5}{tC}{-90}\bike{1.5}{1}{tD}{180}
  \fill[black, opacity=0.55, rounded corners=3pt] (-0.5,-0.5) rectangle (5.5,5.5);
\end{scope}
\begin{scope}[shift={(30.4,0)}]
  \foreach \i in {0,...,5}{\draw[tB, line width=2pt] (\i,0) -- (\i,5); \draw[tB, line width=2pt] (0,\i) -- (5,\i);}
\end{scope}
\node[font=\scriptsize, align=center] at (2.5,-2.4) {stores glow,\\roads dark};
\node[font=\scriptsize, align=center] at (10.1,-2.4) {light moves in,\\3, 2, 1};
\node[font=\scriptsize, align=center] at (17.7,-2.4) {the live game};
\node[font=\scriptsize, align=center] at (25.3,-2.4) {all dims and\\breathes slowly};
\node[font=\scriptsize, align=center] at (32.9,-2.4) {roads in the\\winner colour};
\end{tikzpicture}
\end{ckfig}

In the lobby the roads are dark and the four stores glow in the team colours. In the countdown the
light moves from the edge to the centre in three steps, one for each number. During the game the board
shows the live state. In a pause everything stays where it is, but all LEDs dim and breathe slowly. At
the end all roads turn to the colour of the winner team, and in a draw the colours of the winners take
turns.

\msub{When the link is lost}
If no state comes for a short time, M6 stops the motors at once and the link state becomes
``searching''. The roads go dark and the stores breathe slowly in a dim white, so everybody can see
that the board is waiting. Then M6 tries to connect again: first to the server it finds on the local
network, then to the backup address. The Wi-Fi name, the password, the backup address and the
calibration values are stored in the ESP32's flash memory. We can change them over the USB cable, so we
do not need to build the code again at a new place.

\begin{ckfig}
\begin{tikzpicture}[ck]
\node[st, draw=black!50, fill=black!6] (b) at (0,0) {Start};
\node[st, draw=adm, fill=admL] (w) at (3.4,0) {Join Wi-Fi};
\node[st, draw=adm, fill=admL] (f) at (6.8,0) {Find server};
\node[st, draw=cli, fill=cliL, line width=1.4pt] (c) at (10.2,0) {Connected};
\node[st, draw=obs, fill=obs!10] (s) at (6.8,-1.8) {Searching};
\draw[inna,->] (b) -- (w);
\draw[inna,->] (w) -- node[al]{ok} (f);
\draw[inna,->] (f) -- node[al]{found} (c);
\draw[inna,->] (c) |- node[al, pos=0.3]{no state} (s);
\draw[inna,->] (s) -- node[al]{retry} (f);
\draw[inna,->] (s) -| node[al, pos=0.75]{no Wi-Fi} (w);
\ckleg{0,-2.7}{inna/change of the link state}
\end{tikzpicture}
\end{ckfig}

\msection{Testing}

M6 is tested without the server in two ways. For the link part two PC programs are used. A server
emulator takes the place of M2. It sends test board states, config and vibrate commands with delay and
loss that we can change, and it shows the earthquake events and the seismic signal. A board emulator
takes the place of the real board, so the server side can also be tested without it. We check the
shared example messages, broken packets, a long run, heavy loss, a router restart and a server restart.
For the sensor we hit each of the 9 zones many times and check that hits are found, that the zone is
right or a neighbour zone, and that the motors do not start a false earthquake. The LED part is tested
on the real board with a built in demo scenario. It gives recorded board states as test data and plays
a full game with every animation and every game phase. Test patterns check the order and place of the
LEDs, the 8 player colours (three people must name each colour correctly), the power limit at full
white and the heat in a long run.

\end{document}
